\chapter{The spec workbook}
\label{ch:workbook}

This chapter is the reference for everything you type into the spec workbook: every
sheet, every column and every setting. For each one it says whether you must fill it,
which values are legal, and what the value means. Where a check compares rows with each
other, the text names the rule ID of that check. Chapter~\ref{ch:checks} explains how a
finding reads and lists every rule.

The workbook described here is schema 0.7.

% Column tables of this chapter: column (or setting), required, values, meaning.
\newenvironment{wbtable}[1]{%
  \par\small
  \begin{longtable}{@{}>{\raggedright\arraybackslash}p{3.3cm}%
      >{\raggedright\arraybackslash}p{1.7cm}%
      >{\raggedright\arraybackslash}p{4.1cm}%
      >{\raggedright\arraybackslash}p{5.6cm}@{}}
  \toprule
  \textbf{#1} & \textbf{Required} & \textbf{Values} & \textbf{Meaning}\\
  \midrule
  \endhead
  \bottomrule
  \endfoot}%
  {\end{longtable}\par}

\section{How the workbook is laid out}

\subsection{One sheet per topic}

The \term{spec workbook} is an Excel file (\file{.xlsx}). It describes one subsystem:
its clocks, its bus fabrics, the ports on them, the address map and the wrapper around
them. \cmd{bag template} writes a blank workbook, and \cmd{bag template --sample}
writes the \code{demo_ss} sample of chapter~\ref{ch:first}. Excel and LibreOffice can
open it.

The workbook has one sheet per topic, in this order:

\begin{center}
\small
\begin{tabular}{@{}ll@{}}
\toprule
\textbf{Sheet} & \textbf{What it holds}\\
\midrule
\sheet{README} & A legend and a short guide. BAG does not read it.\\
\sheet{Project} & Project-wide settings and the outputs to generate.\\
\sheet{Clocks} & One row per clock, and so per clock domain.\\
\sheet{Resets} & One row per reset.\\
\sheet{Fabrics} & One row per bus fabric.\\
\sheet{Masters} & One row per external master port.\\
\sheet{Slaves} & One row per external slave port.\\
\sheet{Bridges} & One row per bridge between two fabrics.\\
\sheet{Topology} & One row per link: what attaches to which fabric.\\
\sheet{Access_Rules} & Exceptions to the default access of masters to slaves.\\
\sheet{Memory_Map} & One row per address region.\\
\sheet{Address_Decode} & One row per master port: its miss response and alias.\\
\sheet{Port_Options} & Buffering and optional signals of a port.\\
\sheet{User_Ports} & Extra ports of a fabric that you write yourself.\\
\sheet{Hierarchy} & Sub-wrappers, and which instance goes in each.\\
\sheet{Wrapper} & The top-level wrapper and its exposed ports.\\
\sheet{Lists} & The values of the dropdown lists. BAG does not read it.\\
\bottomrule
\end{tabular}
\end{center}

Some terms are used through the whole chapter:

\begin{description}
\item[Fabric] One bus instance: a crossbar, a shared bus, an AHB-Lite bus or an APB
  decoder.
\item[Master] A port that issues transactions into a fabric, such as a CPU or a DMA
  engine.
\item[Slave] A port that receives transactions from a fabric, such as a memory or a
  peripheral.
\item[Bridge] A link from one fabric to another. It is a slave on its
  \term{upstream} fabric and a master on its \term{downstream} fabric.
\item[Port] A master, a slave, or one end of a bridge.
\item[Clock domain] A group of logic on one clock. Every fabric and every master and
  slave is in exactly one domain.
\item[Adapter] A small block that BAG puts between a port and its fabric: a protocol
  adapter, a width converter, a clock-domain crossing (CDC), an ID-width converter, a
  FIFO or a register slice.
\end{description}

\subsection{Input columns and derived columns}

Each sheet has a title in row~1 and a short description in row~2. Row~4 holds the
column headers, and data starts in row~5. Leave rows 1 to 3 as they are.

There are two kinds of column:

\begin{itemize}
\item An \term{input column} has a blue header and yellow cells. You type into it. Many
  input cells have a dropdown list. The comment on a header gives the meaning and the
  legal values.
\item A \term{derived column} has a grey header and grey cells. It holds a formula that
  shows what BAG will work out, for example the fabric a master is on, or the adapters
  BAG will insert. Do not type into it. BAG ignores derived columns and computes every
  derived value again itself.
\end{itemize}

BAG finds each input column by its header text, not by its position. You can move
columns, and you can add columns of your own: BAG ignores a column whose header it does
not know. A header must be spelled exactly as in this chapter. A missing or repeated
input column stops \cmd{bag check} with a message such as
\code{Fabrics: missing column 'Topology' in row 4 (READ-COLUMN)}. A missing sheet
gives \code{READ-SHEET}.

BAG reads a row when at least one of its input cells has a value. Empty rows, and rows
that only hold derived formulas, are skipped. BAG reads what is typed in a cell, not the
result of a formula: type values, not formulas, into input cells.

The template fills the grey formulas a few rows below the last data row, and its lookups
and dropdowns cover rows 5 to 60 (5 to 200 on \sheet{Topology}). If you add rows
beyond those, copy the formulas down and widen the ranges. BAG reads and checks every
row in any case.

\subsection{Check columns}

Some sheets have a derived \textsf{Check} column. It reads \code{ok} in green, or the name
of the first problem in red, for example \code{OVERLAP}. It is a preview in the
spreadsheet. \cmd{bag check} runs the same checks and many more, and only its result
counts.

\subsection{Notes columns}

Most tables end with a \col{Notes} column, and \sheet{Masters} and \sheet{Slaves} also
have \col{Description}. Type anything you like there. BAG reads these cells but does not
use them, and they are not part of the IR. A workbook that BAG writes back from an IR
(chapter~\ref{ch:ir}) has empty Notes.

\subsection{How to type values}

\begin{description}
\item[Names] Instance, port and region names are lowercase SystemVerilog identifiers: a
  lowercase letter, then lowercase letters, digits and \code{_}. For example
  \code{main_xbar} or \code{cpu0}. A name that does not have this form gives
  \ruleid{NAME-IDENT}. Clock, reset and domain names are SystemVerilog identifiers that
  may also use capitals, such as \code{D_SYS}.
\item[Unique names] The names of masters, slaves, bridges, fabrics and sub-wrappers
  share one name space: each name may be used once (\ruleid{NAME-UNIQUE}).
\item[Hex] Addresses and sizes are hex text with a \code{0x} prefix. You may group the
  digits with \code{_}, for example \code{0x4000_0000}.
\item[Integers] Widths, depths and counts are whole numbers. Widths are in bits and
  frequencies in MHz.
\item[Y and N] Yes or no is a capital \code{Y} or \code{N}.
\item[Blank and n/a] A blank cell means \emph{not given}. Some columns also accept
  \code{n/a} for a value that does not apply, for example \col{Latency mode} on a fabric
  that is not a crossbar.
\end{description}

A required cell that is empty gives \ruleid{VALUE-MISSING}. A value that is not legal
for its column gives \ruleid{VALUE-ILLEGAL}.

\subsection{The order to fill the sheets}

Later sheets refer to names that earlier sheets define, so fill them in this order:

\begin{enumerate}
\item \sheet{Project}: name the subsystem and choose the outputs.
\item \sheet{Clocks} and \sheet{Resets}: every clock domain and its reset.
\item \sheet{Fabrics}, \sheet{Masters}, \sheet{Slaves} and \sheet{Bridges}: define every
  fabric and port by name. Their grey fabric columns read \code{NOT CONNECTED} or
  \code{?} until the next step.
\item \sheet{Topology}: connect everything, one link per row.
\item \sheet{Access_Rules}: only the exceptions, if you have any.
\item \sheet{Memory_Map}: the regions of every slave and bridge.
\item \sheet{Address_Decode}: the miss response of every master port, and any alias.
\item \sheet{Port_Options}: FIFOs, register slices and optional signals, if needed.
\item \sheet{User_Ports}: only for a fabric you write yourself.
\item \sheet{Hierarchy}: sub-wrappers, if you want them.
\item \sheet{Wrapper}: the top module and its ports.
\end{enumerate}

Then run \cmd{bag check} and fix every error before you run \cmd{bag generate}.

\section{README}

The \sheet{README} sheet holds the title of the workbook, a legend of the cell colours
and one line about each sheet. It is for people. BAG does not read it, so you may change
it freely.

\section{Project}

The \sheet{Project} sheet is a list of settings. Each row has a \code{Setting}, a
\code{Value} and a \code{Notes} cell; you type only into \code{Value}. BAG reads the
settings from row~5 down to the first empty \code{Setting} cell. Every setting must be
present (\ruleid{VALUE-MISSING}) and each may appear once.

\begin{wbtable}{Setting}
\setting{Project name} & Yes & lowercase name, e.g.\ \code{demo_ss} &
  The name of the subsystem. It starts the names of several generated files, such as
  \file{<project>.f} and \file{sw/<project>_memmap.h}.\\
\setting{Spec schema version} & Yes & \code{0.7} &
  The schema of the workbook. The template fills it in.\\
\setting{Spec revision} & No & any text, e.g.\ \code{r4} &
  Printed in the header of every generated file and in the memory-map document.\\
\setting{Author} & No & any text &
  Kept in the IR.\\
\setting{SystemVerilog standard} & Yes & \code{IEEE 1800-2017} &
  The language standard of the generated RTL. It is the only legal value.\\
\setting{Output directory} & Yes & a path, e.g.\ \code{out/demo_ss} &
  Where \cmd{bag generate} writes. A relative path starts at the folder of the
  workbook.\\
\setting{Generate RTL} & Yes & \code{Y}, \code{N} &
  Kept in the IR. See below.\\
\setting{Generate C header} & Yes & \code{Y}, \code{N} &
  \code{Y} writes a C header with the base address and size of every region.\\
\setting{Generate memory-map doc} & Yes & \code{Y}, \code{N} &
  \code{Y} writes the memory map (Markdown and HTML), the port map, the tie-offs and the
  reset sequence documents.\\
\setting{Generate cocotb unit tests} & Yes & \code{Y}, \code{N} &
  \code{Y} writes cocotb tests for this design.\\
\setting{Generate formal (SymbiYosys)} & Yes & \code{Y}, \code{N} &
  \code{Y} writes SymbiYosys formal jobs.\\
\setting{Lint tools} & Yes & a comma-separated list of \code{slang}, \code{verilator},
  \code{verible} &
  The tools \cmd{bag lint} runs. Name each tool once.\\
\setting{Simulator} & Yes & \code{verilator} &
  The simulator of \cmd{bag sim}. It is the only legal value.\\
\setting{Module prefix} & No & lowercase name, e.g.\ \code{demo_ss}, without a trailing
  underscore &
  Meant to go, with an underscore, in front of the BAG and PULP module, package and macro
  names and their file names, so that two BAG subsystems or another copy of PULP can be
  compiled together. \emph{This version does not apply it to anything yet}: it only reads
  and checks the value, and the generated files are the same with or without it. Blank is
  legal. \ruleid{PROJ-MODULE-PREFIX} refuses a value that is not a lowercase identifier or
  ends with an underscore. \ruleid{PROJ-PREFIX-COLLISION} refuses a prefix under which a
  sub-wrapper or a BAG or PULP unit would get the name of the top wrapper, the parameter
  package, an interface, a user module or a SystemVerilog keyword; the top wrapper, the
  parameter package and user modules are never renamed. \ruleid{PROJ-NAME-LENGTH} warns
  when a prefixed name would be longer than 64 characters.\\
\end{wbtable}

If \setting{Spec schema version} is not \code{0.7}, \cmd{bag check} reports
\ruleid{PROJ-SCHEMA-VERSION} and checks nothing else: a workbook of another schema has
other columns. Make a new template with \cmd{bag template} and copy your rows into it.

With \setting{Generate RTL} set to \code{N}, \cmd{bag generate} writes only the C header and
the documents that need no RTL. The option \opt{-o} of \cmd{bag generate} replaces
\setting{Output directory} for one run. Chapter~\ref{ch:outputs} lists every file that
the \code{Generate} settings control.

A design with a fabric that you write yourself (template \code{bag_external}) gets no
generated tests and no formal jobs, whatever the settings say: BAG does not know how
that fabric behaves.

\section{Clocks and Resets}

Each clock defines one clock domain, and each domain has exactly one reset. Every clock
and every reset becomes a port of the top-level wrapper, with the name you give it.

\subsection{Clocks}

The \sheet{Clocks} sheet has one row per clock.

\begin{wbtable}{Column}
\col{Clock} & Yes & identifier, e.g.\ \code{clk_sys} & The name of the top-level clock
  port.\\
\col{Freq (MHz)} & Yes & a number above 0 & The nominal frequency. BAG keeps it in the
  IR; it does not use it.\\
\col{Source} & No & any text & Where the clock comes from, for example a PLL.\\
\col{Domain} & Yes & identifier, e.g.\ \code{D_SYS} & The name of the clock domain.
  Other sheets refer to it.\\
\col{Reset} & Yes & a reset of \sheet{Resets} & The reset of this domain.\\
\col{Notes} & No & any text & For you.\\
\end{wbtable}

Each domain has one clock: a second clock for the same domain gives
\ruleid{CLK-DOMAIN-TWO-CLOCKS}. A reset that is not on \sheet{Resets} gives
\ruleid{REF-UNKNOWN}. A reset whose own \col{Domain} names another domain gives
\ruleid{CLK-RESET-DOMAIN}. Clock and reset names must differ from each other
(\ruleid{NAME-UNIQUE}).

\subsection{Resets}

The \sheet{Resets} sheet has one row per reset.

\begin{wbtable}{Column}
\col{Reset} & Yes & identifier, e.g.\ \code{rst_sys_n} & The name of the top-level
  reset port.\\
\col{Polarity} & Yes & \code{active-low}, \code{active-high} & The level that resets.\\
\col{Assertion} & Yes & \code{async}, \code{sync} & How the reset asserts. See
  below.\\
\col{Deassertion} & Yes & \code{async}, \code{sync} & How the reset deasserts. See
  below.\\
\col{Domain} & Yes & a domain of \sheet{Clocks} & The domain this reset belongs to.\\
\col{Source} & No & any text & Where the reset comes from.\\
\col{Notes} & No & any text & For you.\\
\end{wbtable}

Inside the wrapper, every domain gets one active-low reset. With \col{Deassertion}
\code{sync}, BAG puts a two-stage reset synchronizer on the domain's clock. Then
\col{Assertion} \code{async} asserts the reset at once, and \code{sync} asserts it on
the clock, after the two stages. With \col{Deassertion} \code{async}, BAG uses the reset
input as it is (inverted if it is active-high).

Rules across rows: a domain that is not on \sheet{Clocks} gives \ruleid{REF-UNKNOWN}; a
second reset for one domain gives \ruleid{CLK-RESET-DOMAIN}; a domain without a reset
gives \ruleid{CLK-DOMAIN-NO-RESET}.

A clock-domain crossing has logic in both domains, which must be reset together. So the
resets on both sides of every crossing (a CDC adapter, or a bridge between fabrics in
different domains) must have \col{Assertion} \code{async} and \col{Deassertion}
\code{sync} (\ruleid{CLK-CDC-RESET}). That the two resets really assert at the same time
is your job: the workbook cannot show it.

\section{Fabrics}

The \sheet{Fabrics} sheet has one row per fabric. The \term{backing template} names the
BAG template that implements the fabric. You do not type the number of ports: BAG counts
the links of the fabric on \sheet{Topology}.

\begin{wbtable}{Column}
\col{Fabric} & Yes & lowercase name, e.g.\ \code{main_xbar} & The instance name.\\
\col{Protocol} & Yes & \code{AXI4}, \code{AXI4-Lite}, \code{AHB-Lite}, \code{APB4},
  \code{APB3} & The protocol inside the fabric.\\
\col{Topology} & Yes & \code{crossbar}, \code{shared-bus}, \code{single-layer},
  \code{decoder-mux} & The structure of the fabric. See below.\\
\col{Data width (bits)} & Yes & AXI4: 32 to 1024, a power of 2. AXI4-Lite, AHB-Lite: 32
  or 64. APB: 8, 16 or 32 & The data width inside the fabric.\\
\col{Addr width (bits)} & Yes & 12 to 64 & The address width inside the fabric.\\
\col{Clock domain} & Yes & a domain of \sheet{Clocks} & The domain of the fabric.\\
\col{Arbitration} & Yes & \code{round-robin}, \code{fixed-priority},
  \code{single-master} & How the fabric chooses between masters. See below.\\
\col{Latency mode} & Crossbar & \code{NO_LATENCY}, \code{CUT_SLV_AX}, \code{CUT_MST_AX},
  \code{CUT_ALL_AX}, \code{CUT_SLV_PORTS}, \code{CUT_MST_PORTS}, \code{CUT_ALL_PORTS};
  \code{n/a} otherwise & Where the crossbar has spill registers.\\
\col{Max mst trans} & AXI & integer, 1 or more; \code{n/a} otherwise & Limits the
  transactions in flight from each master. See below.\\
\col{Max slv trans} & AXI & integer, 1 or more; \code{n/a} otherwise & Limits the
  transactions in flight towards slaves. See below.\\
\col{Backing template} & Yes & a fabric template, see the table below & The template
  that implements the fabric.\\
\col{Backing IP} & No & any text & For you: the IP behind the template.\\
\col{Mid-point FIFO depth} & Shared bus & integer, 0 or more; \code{n/a} otherwise & The
  FIFO between the mux and the demux of a shared bus.\\
\col{Mid-point register slice} & Shared bus & \code{none}, \code{cut},
  \code{multicut-2}, \code{multicut-3}; \code{n/a} otherwise & Register stages between
  the mux and the demux of a shared bus.\\
\col{Notes} & No & any text & For you.\\
\col{External module} & No & identifier, or blank or \code{n/a} & Only for
  \code{bag_external}: the name of your module. See section~\ref{sec:wb-user-ports}.\\
\end{wbtable}

A setting that does not apply to the fabric must be blank or \code{n/a}. If it holds a
value, \cmd{bag check} warns with \ruleid{FAB-NA-SETTING}: set it to \code{n/a}. This
covers \col{Latency mode} outside a crossbar, \col{Max mst trans} and
\col{Max slv trans} on AHB-Lite and APB, the two mid-point settings outside a shared bus,
and \col{External module} on a template other than \code{bag_external}. Where one of the
other settings applies, it is required.

The grey columns show \textsf{\# master ports} and \textsf{\# slave ports} (the links
into and out of the fabric on \sheet{Topology}), the master-side and slave-side ID
widths of an AXI4 fabric, and \textsf{Max parallel transfers}. These are explained
below.

\subsection{Fabric templates}

Each fabric template supports some protocols, topologies and arbitration modes:

\begin{center}
\small
\begin{tabular}{@{}lll>{\raggedright\arraybackslash}p{4.8cm}@{}}
\toprule
\textbf{Template} & \textbf{Protocol} & \textbf{Topology} & \textbf{Arbitration}\\
\midrule
\code{bag_axi_xbar} & AXI4 & crossbar & round-robin\\
\code{bag_axil_xbar} & AXI4-Lite & crossbar & round-robin\\
\code{bag_axi_shared_bus} & AXI4 & shared-bus & round-robin, fixed-priority\\
\code{bag_axil_shared_bus} & AXI4-Lite & shared-bus & round-robin, fixed-priority\\
\code{bag_ahb_lite_bus} & AHB-Lite & single-layer & single-master, round-robin,
  fixed-priority\\
\code{bag_apb_mux} & APB4, APB3 & decoder-mux & single-master\\
\code{bag_external} & any & any & any\\
\bottomrule
\end{tabular}
\end{center}

The AXI fabrics use the PULP \code{axi} library. The AHB-Lite bus and the APB decoder
are written for BAG. \code{bag_external} is a fabric that you write yourself
(section~\ref{sec:wb-user-ports}).

The rules:

\begin{itemize}
\item The \col{Topology} must suit the \col{Protocol}: \code{crossbar} or
  \code{shared-bus} for AXI4 and AXI4-Lite, \code{single-layer} for AHB-Lite,
  \code{decoder-mux} for APB (\ruleid{FAB-TOPOLOGY}).
\item The template must be a fabric template that supports the protocol and the
  topology (\ruleid{TMPL-FABRIC}), and the arbitration (\ruleid{FAB-ARBITRATION}).
\item \code{single-master} allows one master port. An APB fabric always allows only one
  (\ruleid{FAB-SINGLE-MASTER}).
\item Every fabric needs at least one master port and one slave port on
  \sheet{Topology} (\ruleid{FAB-EMPTY-SIDE}).
\item The data width must be legal for the protocol (\ruleid{WIDTH-DATA}) and the
  address width must be 12 to 64 (\ruleid{WIDTH-ADDR}).
\end{itemize}

\subsection{Crossbar or shared bus}

For AXI4 and AXI4-Lite you choose between two topologies.

A \term{crossbar} has a path from every master to every slave. Transfers between
different masters and slaves run at the same time: up to the smaller of the number of
master ports and slave ports, which the grey column \textsf{Max parallel transfers}
shows. A slow slave stalls only the masters that use it. The logic grows with the number
of masters times the number of slaves. Use a crossbar for the main system interconnect.

A \term{shared bus} takes all masters through one mux into a single mid-point, then
through one demux to the slaves. One transfer moves at a time
(\textsf{Max parallel transfers} is 1). The logic grows with the number of masters plus
the number of slaves, so it is smaller. A slow slave can hold up every master. Use a
shared bus for low-bandwidth configuration and register buses, or where area matters
more than bandwidth. \col{Mid-point FIFO depth} sets a FIFO at the mid-point (0: none),
and \col{Mid-point register slice} adds register stages there, which break the timing
path from the mux to the demux. Both serve every master.

\subsection{Arbitration and latency}

\begin{description}
\item[round-robin] Masters that wait take turns.
\item[fixed-priority] The master with the best \col{Priority} on \sheet{Topology} wins.
  0 is the highest priority.
\item[single-master] The fabric has one master port and needs no arbiter.
\end{description}

\col{Latency mode} is a setting of the PULP crossbar. It places spill registers, each
of which adds a clock of latency and breaks a timing path. In PULP names, the slave
ports of the crossbar face the masters, and its master ports face the slaves.
\code{NO_LATENCY} has no spill registers. \code{CUT_SLV_AX}, \code{CUT_MST_AX} and
\code{CUT_ALL_AX} put them on the address channels (AW and AR) of the slave ports, of
the master ports, or of both. \code{CUT_SLV_PORTS}, \code{CUT_MST_PORTS} and
\code{CUT_ALL_PORTS} put them on all five channels. The sample uses \code{CUT_ALL_AX}.

\col{Max mst trans} limits the transactions in flight from each master:

\begin{itemize}
\item AXI4 crossbar: at least this many per ID, not in total (PULP \code{MaxMstTrans}).
\item AXI4-Lite: this many per master (a FIFO depth).
\item AXI4 shared bus: at least this many per ID at each master, and the depth of the
  write FIFO of the mux, shared by all masters.
\end{itemize}

It also sizes the adapters on master ports.

\col{Max slv trans} limits the transactions towards slaves:

\begin{itemize}
\item AXI4 crossbar: the writes per slave whose data has not finished; reads are not
  limited (PULP \code{MaxSlvTrans}).
\item AXI4-Lite: this many per slave, reads and writes (a FIFO depth).
\item AXI4 shared bus: at least this many per ID, across all slaves.
\end{itemize}

It also limits the reads and writes through a \code{bag_axi2apb} or \code{bag_axi2axil}
bridge on this fabric and the reads through a \code{bag_axi2ahb} bridge, and it sizes
the adapters on slave ports.

\subsection{ID widths}

On an AXI4 fabric, BAG works out two ID widths, shown in grey:

\begin{itemize}
\item The \term{master-side ID width} is the widest \col{ID width (bits)} of the AXI4
  masters on the fabric. A bridge into an AXI4 fabric counts as 1 bit.
\item The \term{slave-side ID width} is the master-side width plus enough bits to
  number the master ports: $\lceil\log_2 n\rceil$ for $n$ master ports.
\end{itemize}

\cmd{bag check} warns when the slave-side width is above 16 bits
(\ruleid{FAB-ID-GROWTH}). On \code{bag_axi_xbar} and \code{bag_axi_shared_bus}, a guard
in front of each master port keeps one counter of transactions in flight per ID, so it
keeps $2^w$ counters per channel for a master-side width $w$. \cmd{bag check} warns when
$w$ is above 8 bits (\ruleid{FAB-GUARD-ID-WIDTH}). That is a cost in area, not a limit.

\section{Masters and Slaves}
\label{sec:wb-ports}

The \sheet{Masters} and \sheet{Slaves} sheets have one row per external port: a port
that leaves the wrapper. The two sheets share most columns. \sheet{Masters} has
\col{Max outstanding rd} and \col{Max outstanding wr}; \sheet{Slaves} has
\col{Error response} and \col{Wait states} in their place.

\begin{wbtable}{Column}
\col{Master}, \col{Slave} & Yes & lowercase name, e.g.\ \code{cpu0} & The port name.
  It also starts the names of the port's wrapper signals.\\
\col{Protocol} & Yes & \code{AXI4}, \code{AXI4-Lite}, \code{AHB-Lite}, \code{APB4},
  \code{APB3} & The protocol of the port.\\
\col{Data width (bits)} & Yes & as on \sheet{Fabrics}, for the port's protocol & The
  data width of the port.\\
\col{Addr width (bits)} & Yes & 1 to 64, not wider than its fabric & The address width
  of the port.\\
\col{ID width (bits)} & AXI4 & integer, 1 or more; on \sheet{Slaves} also \code{auto};
  blank for other protocols & The AXI ID width.\\
\col{User width (bits)} & No & integer, 0 or more; or \code{n/a} & The width of the
  AXI4 USER signals. Blank or 0: none.\\
\col{Clock domain} & Yes & a domain of \sheet{Clocks} & The domain of the port.\\
\col{Max outstanding rd} & AXI (\sheet{Masters}) & integer, 1 or more & Kept in the
  IR.\\
\col{Max outstanding wr} & AXI (\sheet{Masters}) & integer, 1 or more & Kept in the
  IR.\\
\col{Error response} & Yes (\sheet{Slaves}) & AXI: \code{SLVERR}, \code{DECERR}.
  AHB-Lite: \code{HRESP_ERROR}. APB: \code{PSLVERR} & How the slave reports an error.\\
\col{Wait states} & No (\sheet{Slaves}) & integer, 0 or more; or \code{variable} &
  Kept in the IR.\\
\col{Description} & No & any text & For you.\\
\col{Notes} & No & any text & For you.\\
\end{wbtable}

Notes on the columns:

\begin{itemize}
\item The data width must be legal for the port's protocol (\ruleid{WIDTH-DATA}). The
  address width must be 1 to 64 and not wider than the fabric's
  (\ruleid{WIDTH-ADDR}). A narrower slave, such as a 12-bit APB peripheral, sees only
  the low address bits.
\item \col{ID width (bits)} is required on AXI4 ports and must be blank (or
  \code{n/a}) on others (\ruleid{WIDTH-ID}). A slave with \code{auto} gets the
  slave-side ID width of its fabric. A slave with a number that differs from it gets an
  ID-width converter.
\item \col{User width (bits)} above 0 is legal only on AXI4 (\ruleid{WIDTH-USER}).
\item \col{Max outstanding rd}, \col{Max outstanding wr} and \col{Wait states} are
  checked and kept in the IR. The generator does not use them today.
\item \col{Error response} must suit the protocol of the slave
  (\ruleid{VALUE-ILLEGAL}).
\end{itemize}

You do not type the fabric of a port. The grey \textsf{Fabric} column finds it on
\sheet{Topology} and reads \code{NOT CONNECTED} until the port has a link. Every master
and slave needs a link (\ruleid{TOPO-NOT-CONNECTED}).

\subsection{Adapters}

When a port and its fabric differ, BAG puts adapters between them. Adapters exist only
inside the AXI family: between AXI4 or AXI4-Lite ports and AXI4 or AXI4-Lite fabrics.

\begin{itemize}
\item a different clock domain gives a CDC adapter (\code{cdc});
\item a different protocol, AXI4 against AXI4-Lite, gives a protocol adapter
  (\code{pc});
\item a different data width gives a width converter (\code{dw});
\item an AXI4 slave whose ID width differs from the fabric's slave-side ID width gives
  an ID-width converter (\code{iw});
\item a FIFO (\code{fifo}) or a register slice (\code{cut}) comes from
  \sheet{Port_Options}.
\end{itemize}

An AHB-Lite or APB port must have the protocol, data width and clock domain of its
fabric, and an AXI port cannot sit on an AHB-Lite or APB fabric: to cross protocol
families, use a bridge. A mismatch gives \ruleid{PORT-ADAPTER}.

Each adapter has an instance name made from \setting{Adapter instance naming} on
\sheet{Wrapper}: by default \code{<from>_<type>_2_<to>}, for example
\code{cpu0_cdc_2_main_xbar} for a master and \code{main_xbar_dw_2_ddr0} for a slave. Use
these names on \sheet{Hierarchy} to move or rename an adapter. The grey columns on
\sheet{Masters} and \sheet{Slaves} preview the protocol adapter, the width converter and
the CDC with the default pattern. They do not show the ID-width converter, the FIFO or
the register slice, and they do not know the AXI-family rule above. The IR
(chapter~\ref{ch:ir}) lists every adapter BAG inserts.

A width converter cannot split an AXI4 WRAP burst: it answers one with SLVERR. When a
width converter on a master's path refuses its WRAP bursts, \cmd{bag check} warns on
the master's row and names the slaves it concerns (\ruleid{AXI-WRAP-REFUSED}). Chapter~\ref{ch:first} shows the two
warnings of the sample.

\section{Bridges}

The \sheet{Bridges} sheet has one row per bridge. A bridge connects two fabrics. It
appears twice on \sheet{Topology}: once as a destination of its upstream fabric, where
it is a slave, and once as a source into its downstream fabric, where it is a master. So
it needs a region (its \term{window}) on \sheet{Memory_Map}, and a row on
\sheet{Address_Decode} for its downstream end.

\begin{wbtable}{Column}
\col{Bridge} & Yes & lowercase name, e.g.\ \code{br_main2per} & The instance name.\\
\col{Template} & Yes & a bridge template, see the table below & The template that
  implements the bridge.\\
\col{Window base (hex)} & Yes & hex & The lowest base address of the bridge's regions on
  \sheet{Memory_Map}.\\
\col{Pass-through address} & Yes & \code{Y}, \code{N} & \code{Y}: the downstream fabric
  sees the full system address. \code{N}: it sees the offset from the window base.\\
\col{FIFO depth} & No & integer, 0 or more & Requests buffered inside the bridge. Blank
  or 0: none. See below.\\
\col{Notes} & No & any text & For you.\\
\end{wbtable}

The grey columns show the upstream and downstream fabric (from \sheet{Topology}), their
protocols and domains, and whether the bridge crosses clock domains.

\subsection{Bridge templates}

\begin{center}
\small
\begin{tabular}{@{}lllc@{}}
\toprule
\textbf{Template} & \textbf{Upstream fabric} & \textbf{Downstream fabric} &
  \textbf{FIFO depth}\\
\midrule
\code{bag_axi2apb} & AXI4, AXI4-Lite & APB4, APB3 & used\\
\code{bag_axi2ahb} & AXI4, AXI4-Lite & AHB-Lite & used\\
\code{bag_axi2axil} & AXI4 & AXI4-Lite & used\\
\code{bag_ahb2axi} & AHB-Lite & AXI4, AXI4-Lite & ignored\\
\code{bag_ahb2apb} & AHB-Lite & APB4, APB3 & ignored\\
\code{bag_apb2axi} & APB4, APB3 & AXI4, AXI4-Lite & ignored\\
\code{bag_apb2ahb} & APB4, APB3 & AHB-Lite & ignored\\
\bottomrule
\end{tabular}
\end{center}

The template must be a bridge template, and the protocols of the upstream and
downstream fabrics must be ones it supports (\ruleid{TMPL-BRIDGE}).

A bridge does its own width conversion and clock-domain crossing inside, so one bridge
row covers the protocol, the width and the clock. BAG puts no adapters around a bridge,
except a FIFO or register slice from \sheet{Port_Options} on its downstream end. An
error passes back across the bridge in the form of the upstream protocol; for example,
PSLVERR becomes SLVERR on an AXI upstream fabric.

In short, the bridges behave like this:

\begin{itemize}
\item \code{bag_axi2apb} and \code{bag_axi2axil} split bursts into single transfers.
  \code{bag_axi2apb} makes one APB transfer at a time.
\item \code{bag_axi2ahb} takes one AXI transaction at a time (responses in order) and
  maps INCR bursts to AHB INCR bursts.
\item \code{bag_ahb2axi} has one transfer outstanding and issues AXI ID 0 (ID width
  1).
\item \code{bag_ahb2apb}, \code{bag_apb2axi} and \code{bag_apb2ahb} move one transfer
  at a time.
\end{itemize}

The bridges built on PULP blocks refuse some AXI bursts with SLVERR, for example WRAP
bursts. Chapter~\ref{ch:limits} lists these limits.

\subsection{FIFO depth}

\col{FIFO depth} is the depth of a FIFO on every AXI channel inside the bridge. It sits
on the upstream side, after the width conversion and before the clock-domain crossing,
and runs on the upstream clock. With 0 there is no FIFO.

Only \code{bag_axi2apb}, \code{bag_axi2ahb} and \code{bag_axi2axil} have this FIFO. The
other bridges move one transfer at a time through a handshake and ignore the value: a
depth above 0 on them gives the warning \ruleid{BRIDGE-FIFO-IGNORED}.

\subsection{Pass-through address}

\sheet{Memory_Map} always holds full system addresses, on every fabric. The
\col{Pass-through address} setting decides what addresses the downstream fabric decodes.

\begin{itemize}
\item With \code{Y}, the bridge passes the address on as it is. The downstream fabric
  decodes the same system addresses. Its address width must hold them.
\item With \code{N}, the bridge subtracts the window base. The downstream fabric decodes
  offsets from 0, so it can have a small address width. Its regions are still typed as
  system addresses on \sheet{Memory_Map}. They must not lie below the window base, and
  their offsets must fit in the downstream fabric's address width
  (\ruleid{MM-ADDR-RANGE}).
\end{itemize}

A chain of bridges works the same way: a pass-through \code{Y} bridge keeps the offset of
its upstream fabric.

\textbf{One address view per fabric.} A fabric has one decoder, so all of its masters
must send the same kind of address. A fabric fed by a pass-through \code{N} bridge cannot
also have a master that sends other addresses: a direct master, a pass-through \code{Y}
bridge from a fabric without an offset, or a pass-through \code{N} bridge with another
window base. Give them separate fabrics. \cmd{bag check} refuses this at the fabric's row
on \sheet{Fabrics}, naming each master and the addresses it sends
(\ruleid{DEC-MIXED-TRANSLATION}).

\subsection{Rules across rows}

\begin{itemize}
\item Every bridge needs an upstream link and a downstream link
  (\ruleid{TOPO-NOT-CONNECTED}).
\item Bridges must not link fabrics in a loop (\ruleid{TOPO-BRIDGE-CYCLE}). So two
  fabrics are never bridged in both directions.
\item A bridge needs at least one region on \sheet{Memory_Map}, and
  \col{Window base (hex)} must equal the lowest base of its regions
  (\ruleid{BRIDGE-WINDOW-BASE}).
\item Each region of a bridge must reach at least one region on its downstream fabric
  (\ruleid{BRIDGE-WINDOW-EMPTY}).
\end{itemize}

\section{Topology}

The \sheet{Topology} sheet says what attaches where. It has one row per link, read as
\col{Source} to \col{Destination}. It is the only place where you type an attachment:
every grey \col{Fabric} column on other sheets is derived from it.

\begin{wbtable}{Column}
\col{Source} & Yes & a master, fabric or bridge & Where the link starts.\\
\col{Destination} & Yes & a fabric, slave or bridge & Where the link ends.\\
\col{Port index} & Yes & integer, 0 to (ports on that side $-$ 1) & The port number on
  the fabric side of the link.\\
\col{Priority} & Fixed-priority & integer, 0 or more; 0 is highest & The priority of a
  master port on a fixed-priority fabric.\\
\col{Notes} & No & any text & For you.\\
\end{wbtable}

Only four kinds of link are legal:

\begin{center}
\small
\begin{tabular}{@{}ll>{\raggedright\arraybackslash}p{6cm}@{}}
\toprule
\textbf{Link} & \textbf{Example} & \textbf{Meaning}\\
\midrule
master to fabric & \code{cpu0 -> main_xbar} & The master issues into the fabric.\\
fabric to slave & \code{main_xbar -> ddr0} & The fabric serves the slave.\\
fabric to bridge & \code{main_xbar -> br_main2per} & The bridge is a slave on its
  upstream fabric.\\
bridge to fabric & \code{br_main2per -> periph_apb} & The bridge is a master on its
  downstream fabric.\\
\bottomrule
\end{tabular}
\end{center}

A link into a fabric is a \term{master port} of that fabric; a link out of it is a
\term{slave port}. \col{Port index} numbers the master ports of a fabric from 0, and its
slave ports from 0. It sets the port numbers of the RTL and the order of the decode
tables. Because every link has an explicit index, sorting the sheet never renumbers a
port. The port map document lists the ports of each fabric, so a change shows in a
diff.

The rules:

\begin{itemize}
\item Source and Destination must be defined on \sheet{Masters}, \sheet{Slaves},
  \sheet{Bridges} or \sheet{Fabrics} (\ruleid{TOPO-UNKNOWN-NAME}).
\item The link must be one of the four kinds above (\ruleid{TOPO-ILLEGAL-LINK}). A master
  wired straight to a slave is illegal: it needs a fabric between them.
\item A link may appear once (\ruleid{TOPO-DUPLICATE}).
\item Each master, slave and bridge end attaches once (\ruleid{TOPO-MULTI-ATTACH}).
\item \col{Port index} is required (\ruleid{TOPO-INDEX-MISSING}). It must be a whole
  number, 0 or more (\ruleid{TOPO-INDEX-NOT-INTEGER}), below the number of ports on that
  side of the fabric (\ruleid{TOPO-INDEX-OUT-OF-RANGE}), and used once on that side
  (\ruleid{TOPO-INDEX-REPEATED}).
\item A \col{Priority} must be a whole number, 0 or more
  (\ruleid{TOPO-PRIORITY-NOT-INTEGER}). On a fabric with \code{fixed-priority}
  arbitration, every master port needs one (\ruleid{TOPO-PRIORITY-MISSING}), and no two
  master ports of the fabric may share one (\ruleid{TOPO-PRIORITY-REPEATED}). On other
  links, BAG ignores the priority.
\end{itemize}

The grey \textsf{Check} column shows the same problems by name, such as
\code{INDEX OUT OF RANGE}. The grey \textsf{Resolved index} column repeats the
\col{Port index}; a blank index is an error (\ruleid{TOPO-INDEX-MISSING}).

\section{Access\_Rules}

By default every master on a fabric can reach every slave on that fabric, limited only
by the \col{Access} of each region on \sheet{Memory_Map}. The \sheet{Access_Rules}
sheet lists only the exceptions. Leave it empty if you have none.

\begin{wbtable}{Column}
\col{Master} & Yes & a master, or a bridge (its downstream end) & The master side of the
  rule.\\
\col{Slave / bridge} & Yes & a slave, or a bridge (its upstream end) & The slave side.
  It must be on the same fabric as the master.\\
\col{Access} & Yes & \code{RO}, \code{WO}, \code{-} & \code{RO}: reads only.
  \code{WO}: writes only. \code{-}: no access at all.\\
\col{Notes} & No & any text & For you.\\
\end{wbtable}

An access that a rule blocks gets the master's miss response from
\sheet{Address_Decode}. With \code{-}, \code{bag_axi_xbar} also leaves that path out of
the crossbar.

The rules: both names must be connected ports (\ruleid{ACC-UNKNOWN-NAME}) on the same
fabric (\ruleid{ACC-NOT-SAME-FABRIC}), and a pair may have one rule
(\ruleid{ACC-DUPLICATE}). If every master of a fabric is blocked from a slave,
\cmd{bag check} warns (\ruleid{ACC-UNREACHABLE-SLAVE}).

\section{Memory\_Map}

The \sheet{Memory_Map} sheet has one row per address \term{region}. All addresses are
full system addresses, so a peripheral behind a bridge has the same address on every
fabric. A slave may have several regions. A bridge needs at least one: its window on the
upstream fabric.

\begin{wbtable}{Column}
\col{Slave / bridge} & Yes & a slave or bridge & The owner of the region. It must be
  connected on \sheet{Topology}.\\
\col{Region} & Yes & lowercase name, e.g.\ \code{regs} & The region name, unique per
  slave. The C header names it \code{<SLAVE>_<REGION>_BASE}.\\
\col{Base (hex)} & Yes & hex, e.g.\ \code{0x4000_0000} & The first address. It must be
  a multiple of the size.\\
\col{Size (hex)} & Yes & hex, a power of 2, e.g.\ \code{0x0000_1000} & The size in
  bytes.\\
\col{Access} & Yes & \code{RW}, \code{RO}, \code{WO} & The access every master gets.
  \sheet{Access_Rules} can narrow it per master.\\
\col{Secure only} & Yes & \code{Y}, \code{N} & \code{Y}: a non-secure access gets the
  miss response.\\
\col{Notes} & No & any text & For you.\\
\end{wbtable}

You do not type the fabric of a region: it is the fabric of the slave or bridge, from
\sheet{Topology}. It still matters. It decides which decoder holds the region, and it
scopes the overlap check: a bridge window and the slaves behind it share addresses, but
on different fabrics. The grey columns show the end address, the fabric, decimal helpers
and a \textsf{Check}.

The rules for each region:

\begin{itemize}
\item The owner must be on \sheet{Slaves} or \sheet{Bridges} (\ruleid{MM-UNKNOWN-SLAVE})
  and have a link as a destination (\ruleid{MM-NOT-CONNECTED}).
\item A slave may use a region name once (\ruleid{NAME-UNIQUE}).
\item The size must be a power of 2 (\ruleid{MM-SIZE-NOT-POW2}), and the base a
  multiple of the size (\ruleid{MM-MISALIGNED}).
\item The region must fit in the address space of its fabric (after the offset of a
  pass-through \code{N} bridge), and its size must fit in the address width of its
  slave (\ruleid{MM-ADDR-RANGE}).
\item A slave narrower than its fabric sees only its low address bits. So all regions of
  that slave must lie in one aligned block of $2^w$ bytes, for an address width $w$
  (\ruleid{MM-SLAVE-ALIGN}). Otherwise two regions would look the same to the slave.
\item Two regions on one fabric must not overlap (\ruleid{MM-OVERLAP}).
\item On an AXI4 or AXI4-Lite fabric, a region smaller than 4~KB gives a warning
  (\ruleid{MM-SMALL-AXI}): AXI bursts must not cross a 4~KB boundary, and small regions
  make that easy to break.
\item A connected slave without any region gives a warning (\ruleid{MM-PATH-NO-REGION}).
\end{itemize}

\textbf{Secure only.} A port marks an access as non-secure with AxPROT[1] on AXI, HNONSEC
on AHB-Lite and PPROT[1] on APB4. A secure-only region must not be reachable from a
master port that cannot drive that bit: one whose \col{PROT} is tied on
\sheet{Port_Options}, or an APB3 port, which has no PPROT
(\ruleid{SEC-PROT-TIED}). A tied PROT always marks accesses as non-secure.

\section{Address\_Decode}

The \sheet{Address_Decode} sheet has one row for every master port: every master, and
the downstream end of every bridge. It says what a master gets for an address that it
may not use, and it can add one alias.

\begin{wbtable}{Column}
\col{Master} & Yes & a master, or a bridge (its downstream end) & The master port.\\
\col{Miss response} & Yes & \code{DECERR}, \code{SLVERR}, \code{HRESP_ERROR},
  \code{PSLVERR}, \code{default-slave} & What a miss returns. It must suit the fabric's
  protocol.\\
\col{Default slave} & With default slave & a slave on the same fabric & The slave
  that receives every miss.\\
\col{Alias base (hex)} & For an alias & hex & The address the master issues.\\
\col{Alias size (hex)} & For an alias & hex, a power of 2 & The size of the alias
  window.\\
\col{Alias maps to (hex)} & For an alias & hex & The address decoded instead.\\
\col{Alias enable} & For an alias & \code{always}, or an extra top-level input & When the
  alias is active.\\
\col{Notes} & No & any text & For you.\\
\end{wbtable}

The fabric is derived and shown in grey.

\subsection{Miss response}

A \term{miss} is an access that no region accepts: the address is outside every region
of the fabric, or an access rule or the region's \col{Access} blocks it, or it is
non-secure and the region is secure only. The miss response must suit the protocol of
the master's fabric (\ruleid{DEC-MISS-RESP}):

\begin{itemize}
\item AXI4 and AXI4-Lite: \code{DECERR}, \code{SLVERR} or \code{default-slave};
\item AHB-Lite: \code{HRESP_ERROR} or \code{default-slave};
\item APB: \code{PSLVERR} or \code{default-slave}.
\end{itemize}

For a bridge, it is the protocol of the downstream fabric. The bridge carries the error
back to its upstream master in the upstream protocol.

With an error response, the fabric answers the miss itself. With \code{default-slave},
the miss goes to the slave you name in \col{Default slave} instead, and that slave
answers. Use it when a real slave should see stray accesses, for example to log them.
The default slave must be a slave on the same fabric, and \col{Default slave} must be
blank with any other miss response (\ruleid{DEC-DEFAULT-SLAVE}).

Every master port needs exactly one row (\ruleid{DEC-MISSING-ROW},
\ruleid{DEC-DUPLICATE}), and the name must be a connected master or bridge
(\ruleid{DEC-UNKNOWN-MASTER}).

\subsection{Aliases}

An \term{alias} makes a master see a window of addresses as another window. For example,
in the sample, \code{cpu0} sees \code{sram0} (at \code{0x1000_0000}) also at
\code{0x0000_0000} while the input \code{boot_remap_i} is 1. An alias applies only to
the master of its row; other masters do not see it. A master can have one alias.

Fill all four alias columns, or none (\ruleid{VALUE-MISSING}). The rules:

\begin{itemize}
\item The size must be a power of 2, and the alias base and the target must be
  multiples of it (\ruleid{DEC-ALIAS-SIZE}).
\item The alias window must not overlap a region on the fabric
  (\ruleid{DEC-ALIAS-OVERLAP}).
\item The target window must lie inside one region on the fabric
  (\ruleid{DEC-ALIAS-TARGET}).
\item \col{Alias enable} is \code{always}, or the name of an input listed in
  \setting{Extra top-level inputs} on \sheet{Wrapper} (\ruleid{DEC-ALIAS-ENABLE}).
\end{itemize}

An \term{alias pin} is an enable input of the wrapper. Use one when software or a boot
strap must turn the alias on and off, for example a boot remap that ends after reset.
Use \code{always} for an alias that is fixed.

\section{Port\_Options}

The \sheet{Port_Options} sheet changes the defaults of a port. List only the ports that
need it. A port that is not listed gets no FIFO, no register slice, and every optional
signal of its protocol on the wrapper.

\begin{wbtable}{Column}
\col{Port} & Yes & a master, slave or bridge & The port. For a bridge, the options apply
  to its downstream end.\\
\col{FIFO depth} & No & integer, 0 or more & A FIFO on every channel of the port. Blank
  or 0: none.\\
\col{Register slice} & No & \code{none}, \code{cut}, \code{multicut-2},
  \code{multicut-3}, \code{n/a} & Register stages on the port: 1, 2 or 3. Each stage
  adds a clock of latency.\\
\col{PROT} & No & see below & AxPROT (AXI), HPROT and HNONSEC (AHB-Lite), PPROT (APB4).\\
\col{QOS} & No & see below & AxQOS (AXI4).\\
\col{REGION} & No & see below & AxREGION (AXI4).\\
\col{CACHE} & No & see below & AxCACHE (AXI4).\\
\col{LOCK / excl} & No & see below & AxLOCK (AXI4), HMASTLOCK (AHB-Lite masters).\\
\col{PSTRB} & No & see below & PSTRB (APB4).\\
\col{PSLVERR} & No & see below & PSLVERR (APB3, APB4).\\
\col{Notes} & No & any text & For you.\\
\end{wbtable}

\col{FIFO depth} and \col{Register slice} add the adapters \code{fifo} and \code{cut}
between the port and its fabric. They exist only on AXI4 and AXI4-Lite ports
(\ruleid{PORT-ADAPTER}). \col{FIFO depth} adds buffering. It is not the same as
\col{Max mst trans} on \sheet{Fabrics}, which only limits the transactions in flight.

Each optional-signal column takes one of four values. Blank means \code{drive} for a
signal the port has, and \code{n/a} for one it does not have.

\begin{description}
\item[drive] The signal exists and is a port of the wrapper.
\item[tie-0] The signal is not a port of the wrapper; BAG drives 0 inside.
\item[tie-default] The signal is not a port of the wrapper; BAG drives the AMBA default
  inside: AxCACHE and HPROT \code{0b0011}, and PSTRB all ones on writes and 0 on reads.
  Other signals get 0.
\item[n/a] The signal does not exist in the port's protocol.
\end{description}

A tied security bit fails closed: with \code{tie-0} or \code{tie-default} on
\col{PROT}, AxPROT[1] and PPROT[1] are 1 and HNONSEC is 1, so every access is
non-secure.

The rules: the port must be a connected master, slave or bridge
(\ruleid{PORTOPT-UNKNOWN-PORT}), listed once (\ruleid{PORTOPT-DUPLICATE}). Use
\code{n/a} exactly for the signals that the port's protocol lacks, and a real value for
the others (\ruleid{PORTOPT-SIGNAL}). HMASTLOCK goes to AHB-Lite masters only, so
\col{LOCK / excl} is \code{n/a} on an AHB-Lite slave. A tied PSLVERR on a master port
of an APB fabric gives a warning (\ruleid{PORTOPT-ERROR-MASKED}): no error reaches that
master, its own miss response included.

\section{User\_Ports}
\label{sec:wb-user-ports}

A fabric with the backing template \code{bag_external} is one that you write yourself.
BAG puts your module into the wrapper where a generated fabric would be, and connects
its bus ports. It also writes a \term{stub} of the module, \file{rtl/ext/<module>.sv},
with the same ports and every output tied to 0, so that the design lints and elaborates.
Replace the stub with your RTL. BAG writes no RTL, no tests and no formal jobs for such
a fabric; it still writes the memory map and the decode tables.

The module and its bus ports are named like this:

\begin{itemize}
\item With \col{External module} blank, the module is \code{<project>_<fabric>}. Its
  bus ports are named by \setting{Signal naming pattern} with \code{<fabric>_<port>} as
  the port name, for example \code{noc_cpu0_axi_awaddr}.
\item With \col{External module} set, the module has that name, and its bus ports are
  named by position: \code{s<k>} for the k-th master port of the fabric and \code{m<k>}
  for the k-th slave port, for example \code{s0_axi_awaddr}. So fabrics of the same
  shape can share one module.
\item Every such module also has \code{clk_i} and \code{rst_ni}: the fabric's clock and
  its active-low reset.
\end{itemize}

\col{External module} must be a SystemVerilog identifier and not a keyword, and must not
clash with a module BAG generates, start with \code{bag_}, or differ from another
external module only in case (\ruleid{EXT-MODULE-NAME}). Fabrics that share a module
must have the same protocol, widths, number of attachments and user ports
(\ruleid{EXT-MODULE-SHARED}).

The \sheet{User_Ports} sheet declares the other ports of your module: interrupts,
status outputs, configuration inputs, a test mode, extra clocks. It has one row per
port.

\begin{wbtable}{Column}
\col{Instance} & Yes & a fabric with template \code{bag_external} & The fabric whose
  module has the port.\\
\col{Port} & Yes & lowercase name & The port name on your module.\\
\col{Direction} & Yes & \code{input}, \code{output} & As your module sees it.\\
\col{Width (bits)} & Yes & integer, 1 or more & The width of the port.\\
\col{Connection} & Yes & see below & How BAG connects the port.\\
\col{Notes} & No & any text & For you.\\
\end{wbtable}

\begin{center}
\small
\begin{tabular}{@{}ll>{\raggedright\arraybackslash}p{6.5cm}@{}}
\toprule
\textbf{Connection} & \textbf{Direction} & \textbf{Meaning}\\
\midrule
\code{expose} & input or output & A port \code{<instance>_<port>} of the top wrapper.\\
\code{extra input <name>} & input & An input from \setting{Extra top-level inputs}.\\
\code{tie-0}, \code{tie-1}, \code{0x<hex>} & input & A constant. \code{tie-1} sets every
  bit.\\
\code{clock <domain>} & input & The clock of the domain.\\
\code{reset <domain>} & input & The active-low reset of the domain, as BAG's own blocks
  get it.\\
\code{open} & output & Left unconnected.\\
\bottomrule
\end{tabular}
\end{center}

The rules: the instance must be a \code{bag_external} fabric
(\ruleid{USERPORT-INSTANCE}); a port may be declared once per instance
(\ruleid{USERPORT-DUPLICATE}) and must not be one of the bus, clock or reset ports that
BAG already connects (\ruleid{USERPORT-NAME}). The connection must be one of the forms
above and legal for the direction; an extra input or a domain must exist; a clock, a
reset or an extra input is 1 bit wide; and a constant must fit in the width
(\ruleid{USERPORT-CONNECTION}). An exposed name must not clash with another top-level
port (\ruleid{WRAP-NAME-COLLISION}).

A port of your module that is a whole bus, such as an APB configuration port, is
declared here signal by signal.

\section{Hierarchy}

\term{Sub-wrappers} are modules between the top wrapper and the instances. They group
fabrics, bridges and adapters, for example for a power domain or a floorplan partition.
The \sheet{Hierarchy} sheet has two tables: the sub-wrappers, then the placement of
instances in them. You need neither: without them, everything is in the top wrapper.

\subsection{Sub-wrappers}

\begin{wbtable}{Column}
\col{Sub-wrapper} & Yes & lowercase name, e.g.\ \code{main_wrapper} & The module name.
  Row~5 is the top wrapper; its name comes from \setting{Top module name}.\\
\col{Parent} & Yes & \code{-} for the top; a sub-wrapper otherwise & The module that
  contains this one. Sub-wrappers can nest.\\
\col{Notes} & No & any text & For you.\\
\end{wbtable}

The template has rows 5 to 12 for this table. The table ends at the label
\code{Placement} in column A.

The rules: only row~5 has the parent \code{-} (\ruleid{HIER-TOP}); a parent must be a
sub-wrapper or the top (\ruleid{HIER-UNKNOWN-PARENT}), not the sub-wrapper itself
(\ruleid{HIER-SELF-PARENT}); and parents must not form a loop
(\ruleid{HIER-PARENT-LOOP}). A sub-wrapper name must not be the name of another module
of the design (\ruleid{HIER-RESERVED-NAME}). A sub-wrapper that contains no instance,
directly or below it, gives a warning (\ruleid{HIER-EMPTY}).

\subsection{Placement}

The placement table starts below the label \code{Placement}, with its own header row.

\begin{wbtable}{Column}
\col{Instance} & Yes & a fabric, bridge or adapter instance & The instance to place.\\
\col{Sub-wrapper} & Yes & a sub-wrapper, or the top & Where the instance goes.\\
\col{RTL instance name} & No & lowercase name & The instance name in the RTL. Blank:
  \code{u_<instance>}.\\
\col{Notes} & No & any text & For you.\\
\end{wbtable}

List only what you want to move or rename. The defaults are:

\begin{itemize}
\item Fabrics and bridges go in the top wrapper.
\item Adapters follow \setting{Adapter placement} on \sheet{Wrapper}: with
  \code{with fabric}, an adapter goes where its fabric is.
\end{itemize}

Adapter names are those of section~\ref{sec:wb-ports}, for example
\code{main_xbar_dw_2_ddr0}. Where a link crosses a sub-wrapper boundary, BAG makes
ports named by \setting{Sub-wrapper port naming}. An adapter or a bridge is one
instance, so a clock-domain crossing is always generated whole inside one wrapper.

The rules: the instance must exist (\ruleid{HIER-UNKNOWN-INSTANCE}) and be placed once
(\ruleid{HIER-PLACED-TWICE}); the sub-wrapper must be in the sub-wrapper table
(\ruleid{HIER-UNKNOWN-SUBWRAPPER}); and instance names must be unique in each wrapper
(\ruleid{HIER-RTL-NAME}).

\section{Wrapper}

The \sheet{Wrapper} sheet describes the top module that your SoC instantiates. It has a
list of settings, like \sheet{Project}, and then the \code{Exposed ports} table.

\subsection{Settings}

\begin{wbtable}{Setting}
\setting{Top module name} & Yes & lowercase name, e.g.\ \code{demo_ss_bus_top} & The
  top module. Row~5 of \sheet{Hierarchy} shows it.\\
\setting{Port style} & Yes & \code{flat}, \code{sv-interface} & \code{flat}: one port per
  signal. \code{sv-interface}: one SystemVerilog interface per bus port.\\
\setting{Signal naming pattern} & Yes & a pattern with \code{{port}}, \code{{proto}},
  \code{{signal}} & The names of the flat ports, e.g.\ \code{cpu0_axi_awaddr}.\\
\setting{Clock/reset ports} & Yes & \code{one per domain} & One clock and one reset port
  per domain, named as on \sheet{Clocks} and \sheet{Resets}.\\
\setting{Parameter package} & Yes & lowercase name, e.g.\ \code{demo_ss_bus_pkg} & The
  package with the widths, base addresses and region sizes as localparams.\\
\setting{File header} & No & a file name, e.g.\ \code{header.txt} & A text file, such as
  a licence, that starts each generated file.\\
\setting{Extra top-level inputs} & No & comma-separated identifiers, e.g.\
  \code{boot_remap_i} & 1-bit inputs of the top wrapper, for alias enables and user
  ports.\\
\setting{Adapter placement} & Yes & \code{with fabric}, \code{with port}, \code{top} &
  The default wrapper of the adapters.\\
\setting{Adapter instance naming} & Yes & a pattern with \code{{from}}, \code{{type}},
  \code{{to}} & The names of the adapter instances.\\
\setting{Sub-wrapper port naming} & Yes & a pattern with \code{{instance}},
  \code{{proto}}, \code{{signal}} & The names of the ports where a link crosses a
  sub-wrapper boundary.\\
\end{wbtable}

The template fills in the defaults: \code{flat}, \code{{port}_{proto}_{signal}},
\code{one per domain}, \code{with fabric}, \code{{from}_{type}_2_{to}} and
\code{{instance}_{proto}_{signal}}.

\begin{description}
\item[Patterns] A pattern must use each of its placeholders exactly once, and no other
  (\ruleid{WRAP-PATTERN}). In \setting{Signal naming pattern}, \code{{port}} is the port
  name, \code{{proto}} is \code{axi}, \code{axil}, \code{ahb}, \code{apb4} or
  \code{apb3}, and \code{{signal}} is the lowercase AMBA signal name. In
  \setting{Adapter instance naming}, \code{{type}} is \code{pc}, \code{dw}, \code{cdc},
  \code{iw}, \code{fifo} or \code{cut}. In \setting{Sub-wrapper port naming},
  \code{{instance}} is the name of the link: an exposed port's name, or
  \code{<from>_2_<to>} for a link between two instances.
\item[Port style] \code{flat} works with every tool. \code{sv-interface} is more
  compact. With it, BAG also writes one interface per protocol and a flat module around
  the wrapper, \code{<top>_flat}, which lint and compile use as the top. The option
  \opt{--style} of \cmd{bag generate} replaces this setting for one run.
\item[File header] The path starts at the folder of the workbook, and the file must
  exist (\ruleid{WRAP-FILE-HEADER}). Write it as SystemVerilog line comments (\code{//}): BAG copies it as it is to the top of
  the RTL and the C header, and no check reads its text. Documents get it inside an HTML
  comment, tests and formal jobs with \texttt{\#} before each line
  (section~\ref{sec:header}). Copies of library RTL keep their own headers.
\item[Extra top-level inputs] Each name must be a SystemVerilog identifier
  (\ruleid{NAME-IDENT}). The inputs are used by \col{Alias enable} on
  \sheet{Address_Decode} and by \code{extra input} on \sheet{User_Ports}.
\item[Adapter placement] \code{with fabric} puts an adapter in the same wrapper as its
  fabric. \code{with port} puts it beside the port it serves: in the top wrapper for a
  master or slave, which are exposed there, and in the bridge's sub-wrapper for an adapter
  on a bridge end. \code{top} puts it in the top wrapper. \sheet{Hierarchy} can place any adapter elsewhere.
\end{description}

\subsection{Exposed ports}

The \code{Exposed ports} table starts below the settings, with its own header row. It
lists every master, slave and bridge, one per row.

\begin{wbtable}{Column}
\col{Port} & Yes & a master, slave or bridge & The port.\\
\col{Exposed at top} & Yes & \code{Y}, \code{N} & \code{Y} for masters and slaves;
  \code{N} for bridges.\\
\col{Top-level prefix} & No & lowercase name, e.g.\ \code{cpu0_axi} & Replaces
  \code{{port}_{proto}} in the signal naming pattern for this port.\\
\col{If not exposed} & No & \code{internal}, or blank & For you. Only bridges are
  internal.\\
\col{Notes} & No & any text & For you.\\
\end{wbtable}

In this release every master and slave is a port of the top wrapper, and every bridge is
internal (\ruleid{WRAP-INTERNAL-PORT}). A \col{Top-level prefix} takes effect only when
the pattern contains \code{{port}_{proto}}: then \code{cpu0_axi} gives
\code{cpu0_axi_awaddr}, whatever the protocol token.

The rules: every master, slave and bridge must be listed (\ruleid{WRAP-PORT-NOT-LISTED}),
once (\ruleid{WRAP-DUPLICATE-PORT}), and only those (\ruleid{WRAP-UNKNOWN-PORT}). No two
ports of the top wrapper may get the same name, counting clocks, resets, extra inputs
and exposed user ports (\ruleid{WRAP-NAME-COLLISION}). A signal that
\sheet{Port_Options} ties is not a port of the wrapper.

\section{Lists}

The \sheet{Lists} sheet holds the values of the dropdown lists, one list per column:
\code{Protocol}, \code{Topology}, \code{Arbitration}, \code{LatencyMode},
\code{Polarity}, \code{Assertion}, \code{Access}, \code{RuleAccess}, \code{YesNo},
\code{MissResp}, \code{RegSlice}, \code{PortStyle}, \code{Template}, \code{SigDrive},
\code{AdapterPlace} and \code{Direction}. The input cells of the other sheets take their
dropdowns from here.

BAG does not read this sheet. You can add a value below a list and widen the range of the
dropdown, but \cmd{bag check} still refuses a value that BAG does not support. The legal
values are the ones in this chapter.
